\documentclass[12pt]{article}

\usepackage[margin=1in]{geometry}
\usepackage{array}
\usepackage{booktabs}
\usepackage{enumitem}
\usepackage{graphicx}
\usepackage[hidelinks]{hyperref}
\usepackage{parskip}
\usepackage{m\usepackage{longtable}
\usepackage{tabularx}
\usepackage{float}
\usepackage{tikz}
\usetikzlibrary{arrows.meta, positioning, shapes.misc, fit}

\title{\textbf{PodcastHub End-to-End Technical Report}\\Graduate Capstone in Microservice-Oriented Architecture}
\author{Student: Student: \underline{400646346}\\Microservice-Oriented Architecture}
\date{\today}

\begin{document}

\maketitle
\tableofcontents
\listoftables
\listoffigures
\newpage

\section{Executive Overview}

PodcastHub is a distributed podcast studio that demonstrates microservice-oriented, domain-driven, and event-driven architecture in a realistic media workflow. Remote hosts and guests record locally through WebRTC, upload five-second chunks to MinIO, and rely on background workers to stitch final masters using FFmpeg. The system's primary learning objective is to show how clean architectural boundaries enable independent evolution, resilience, and operational clarity.

\section{Architectural Principles}

\begin{table}[H]
\centering
\resizebox{\textwidth}{!}{%
\begin{tabular}{>{\bfseries}p{3.6cm}p{5.8cm}p{6.4cm}}
\toprule
Principle & Implementation Vector & Evidence in Repository \\
\midrule
Microservices &
FastAPI-based Recording API, Processing API, dedicated background worker, and a Next.js frontend deployed as separate containers &
Dockerfiles per service; \texttt{docker-compose.yml} orchestrates services with explicit dependencies and health checks. \\
Domain-Driven Design &
Aggregates (\texttt{Recording}, \texttt{Chunk}, \texttt{Upload}) own invariants while use cases enforce orchestration rules &
\texttt{media-recording-service/src/domain/models}, \texttt{.../application/services}. \\
Hexagonal Architecture &
Ports (\texttt{StoragePort}, \texttt{EventPublisherPort}) decouple infrastructure; adapters implement MinIO, RabbitMQ, repositories &
\texttt{media-recording-service/src/application/ports}, \texttt{.../adapters/outbound}. \\
Event-Driven Coordination &
RabbitMQ topic exchange and command queue separate recording ingest from FFmpeg processing &
\texttt{RabbitMQEventPublisher}, \texttt{processing\_queue.py}, worker consumer. \\
Resilient Storage &
Object storage for binaries, relational database for metadata, durable queues for commands/events &
MinIO integration, PostgreSQL-ready repositories, durable queue definitions in compose file. \\
Container-Native Deployment &
All components packaged with required tooling (Python 3.11, FFmpeg, Node 18) for reproducible demos &
Dockerfiles in each service, environment overrides in \texttt{.env} files. \\
\bottomrule
\end{tabular}}
\caption{Graduate-level architectural principles mapped to implementation artifacts.}
\label{tab:principles}
\end{table}

\section{System Context}

\begin{figure}[H]
\centering
\resizebox{0.92\textwidth}{!}{%
\begin{tikzpicture}[
    node distance=1.8cm,
    service/.style={rectangle, rounded corners, draw=black, fill=blue!10, minimum width=3.2cm, minimum height=1.1cm, align=center},
    infra/.style={rectangle, draw=black, fill=orange!15, minimum width=3.0cm, minimum height=1cm, align=center},
    actor/.style={rectangle, draw=black, fill=green!15, minimum width=3.0cm, minimum height=1cm, align=center},
    arrow/.style={-{Latex[length=3mm]}, thick}
]
    \node[actor] (host) {Host\\Browser};
    \node[actor, below=1.2cm of host] (guest) {Guest\\Browser};

    \node[service, right=3.8cm of host] (recording) {Media Recording\\Service (FastAPI)};
    \node[service, below=1.2cm of recording] (processing) {Media Processing\\Service (FastAPI)};
    \node[service, below=1.2cm of processing] (worker) {Media Processing\\Worker (FFmpeg)};

    \node[infra, right=4.2cm of processing] (rabbit) {RabbitMQ};
    \node[infra, above=1.4cm of rabbit] (minio) {MinIO Object Store};
    \node[infra, below=1.4cm of rabbit] (postgres) {PostgreSQL};
    \node[infra, below=1.0cm of postgres] (redis) {Redis};

    \draw[arrow] (host) -- node[above, sloped]{REST + WebSocket} (recording);
    \draw[arrow] (guest) -- node[below, sloped]{REST + WebSocket} (recording);

    \draw[arrow] (recording) -- node[above]{Chunks} (minio);
    \draw[arrow] (recording) -- node[above]{Events} (rabbit);
    \draw[arrow] (recording) -- node[below]{Metadata} (postgres);

    \draw[arrow] (processing) -- node[above]{Admin REST} (recording);
    \draw[arrow] (processing) -- node[sloped, above]{Monitoring} (rabbit);

    \draw[arrow] (rabbit) -- node[right]{Commands} (worker);
    \draw[arrow] (worker) -- node[left]{Processed Events} (rabbit);
    \draw[arrow] (worker) -- node[below]{Stitched Files} (minio);
\end{tikzpicture}}
\caption{System context: interactions among actors, services, and shared infrastructure.}
\label{fig:context}
\end{figure}

\section{Domain Decomposition}

The platform comprises two bounded contexts: Recording and Processing. The recording context covers live session management and chunk ingestion; the processing context is responsible for offline stitching and downstream orchestration.

\begin{table}[H]
\centering
\resizebox{\textwidth}{!}{%
\begin{tabular}{>{\bfseries}p{3.2cm}p{4.8cm}p{5.6cm}p{6.0cm}}
\toprule
Bounded Context & Ubiquitous Language & Aggregates / Invariants & Primary Use Cases \\
\midrule
Recording &
Session, recording, chunk, upload &
\texttt{Recording}, \texttt{Chunk}, \texttt{Upload} enforce lifecycle transitions, checksum validation, retry ceilings &
Create/join sessions, start/stop recordings, ingest chunks, track progress, enqueue processing commands \\
Processing &
Processing job, stitched asset, workflow &
\texttt{RecordingProcessed} event and future \texttt{ProcessingJob} aggregate preserve provenance and output metadata &
Consume commands, assemble media using FFmpeg, publish completion events, expose processed assets \\
\bottomrule
\end{tabular}}
\caption{Domain-driven decomposition of PodcastHub.}
\label{tab:domain}
\end{table}

\section{Microservice Summary}

\begin{table}[H]
\centering
\resizebox{\textwidth}{!}{%
\begin{tabular}{>{\bfseries}p{3.4cm}p{4.8cm}p{6.0cm}p{5.6cm}}
\toprule
Service & Core Responsibilities & Inbound Ports / Adapters & Outbound Ports / Adapters \\
\midrule
Media Recording Service &
Session orchestration, chunk ingest, progress tracking, command enqueue &
FastAPI REST controllers, WebSocket signalling handler &
MinIO StoragePort adapter, RabbitMQ EventPublisherPort, repository ports \\
Media Processing Worker &
Deterministic FFmpeg stitching, processed upload, event emission &
RabbitMQ consumer (\texttt{media.processing.requests}) &
MinIO StoragePort adapter, RabbitMQ EventPublisherPort \\
Media Processing Service &
Administrative APIs, orchestration entry points, health monitoring &
FastAPI REST controllers &
RabbitMQ command publisher, MinIO client (future workflows) \\
Podcast Frontend (Next.js) &
Host/guest journeys, WebRTC capture, background upload loop, progress UX &
Next.js routes, browser WebRTC APIs, REST client &
Delegates to Recording Service; consumes processed events via polling/webhooks \\
\bottomrule
\end{tabular}}
\caption{Microservice landscape and their hexagonal boundaries.}
\label{tab:services}
\end{table}

\section{Hexagonal Layers}

\begin{figure}[H]
\centering
\resizebox{0.75\textwidth}{!}{%
\begin{tikzpicture}[
    layer/.style={rectangle, draw=black, rounded corners, fill=blue!10, minimum width=9cm, minimum height=1.1cm, align=center},
    adapter/.style={rectangle, draw=black, rounded corners, fill=orange!20, minimum width=3.6cm, minimum height=0.9cm, align=center},
    arrow/.style={-{Latex[length=3mm]}, thick}
]
    \node[layer] (domain) {Domain Aggregates (\texttt{Recording}, \texttt{Chunk}, \texttt{Upload})};
    \node[layer, above=1.4cm of domain] (app) {Application Services (\texttt{RecordingService}, \texttt{UploadService})};

    \node[adapter, above left=1.4cm and 1.9cm of app] (rest) {REST Controllers};
    \node[adapter, above right=1.4cm and 1.9cm of app] (ws) {WebSocket Signalling};
    \node[adapter, below=1.0cm of rest] (cmd) {Processing Command Enqueue};

    \node[adapter, below left=1.4cm and 1.9cm of domain] (storage) {StoragePort (MinIO)};
    \node[adapter, below=1.4cm of domain] (event) {EventPublisherPort (RabbitMQ)};
    \node[adapter, below right=1.4cm and 1.9cm of domain] (repo) {Repository Ports};

    \draw[arrow] (rest) -- (app);
    \draw[arrow] (ws) -- (app);
    \draw[arrow] (cmd) -- (app);
    \draw[arrow] (app) -- (domain);
    \draw[arrow] (domain) -- (storage);
    \draw[arrow] (domain) -- (event);
    \draw[arrow] (domain) -- (repo);
\end{tikzpicture}}
\caption{Hexagonal layering of the Media Recording Service.}
\label{fig:hexagon}
\end{figure}

\clearpage
\section{Microservice Implementation Details}

\subsection{Media Recording Service}
Located in \texttt{media-recording-service/}, this FastAPI application handles session creation, recording lifecycle, chunk uploads, and progress reporting. Domain models live under \texttt{src/domain/models}, use cases under \texttt{src/application/services}, and inbound adapters under \texttt{src/adapters/inbound}. Outbound adapters provide MinIO storage (\texttt{minio\_storage.py}) and RabbitMQ publishing. A simplified HTTP upload route is provided for MinIO demos alongside full REST endpoints.

\subsection{Media Processing Worker}
Packaged within the same codebase, the worker listens to \texttt{media.processing.requests}, downloads chunk objects, builds a concat manifest, invokes FFmpeg (\texttt{ffmpeg -f concat -safe 0 -c copy}), uploads the stitched file to MinIO, and publishes \texttt{recording.processed}. The reuse of domain models underscores the benefit of clean ports/adapters.

\subsection{Media Processing Service}
This FastAPI façade (directory \texttt{media-processing-service/}) exposes administrative endpoints for job creation, manual restarts, status checks, and health monitoring. While the worker currently performs the heavy lifting, the service provides an entry point for future workflows such as mastering or multi-track mixing.

\subsection{Podcast Frontend (Next.js 14)}
Found under \texttt{podcast-frontend/}, the frontend coordinates host/guest flows, WebRTC peer connections, MediaRecorder chunk scheduling, resilient uploads, and progress feedback. Environment variables (\texttt{NEXT\_PUBLIC\_*}) allow targeting local or containerized APIs.

\clearpage
\section{Service Interaction Patterns}

\subsection{Session Lifecycle Narrative}
When a host creates a session, the frontend calls \texttt{POST /api/sessions/create}. The recording service instantiates a \texttt{Recording} aggregate, persists metadata (PostgreSQL-ready), emits \texttt{session.created}, and returns a room code. Guests join via \texttt{/api/sessions/join}, triggering WebSocket signalling that establishes WebRTC peer connections.

\subsection{Recording and Processing Flow}

\begin{figure}[H]
\centering
\resizebox{0.85\textwidth}{!}{%
\begin{tikzpicture}[
    milestone/.style={rectangle, rounded corners, draw=black, fill=gray!10, minimum width=2.8cm, minimum height=1cm, align=center},
    arrow/.style={-{Latex[length=3mm]}, thick}
]
    \node[milestone] (upload) {Chunk Upload};
    \node[milestone, right=2.6cm of upload] (enqueue) {Command Enqueue};
    \node[milestone, right=2.6cm of enqueue] (stitch) {Worker Stitch};
    \node[milestone, right=2.6cm of stitch] (event) {Processed Event};
    \node[milestone, right=2.6cm of event] (delivery) {Download Ready};

    \draw[arrow] (upload) -- node[above]{REST + MinIO} (enqueue);
    \draw[arrow] (enqueue) -- node[above]{RabbitMQ} (stitch);
    \draw[arrow] (stitch) -- node[above]{FFmpeg + MinIO} (event);
    \draw[arrow] (event) -- node[above]{Polling / Event} (delivery);
\end{tikzpicture}}
\caption{End-to-end flow from live chunk upload to processed asset delivery.}
\label{fig:timeline}
\end{figure}

\clearpage
\section{Data Management Strategy}

\begin{itemize}[leftmargin=*]
    \item \textbf{PostgreSQL} (future integration) will hold session, recording, and upload metadata with ACID guarantees.
    \item \textbf{MinIO} stores immutable chunk objects and processed artefacts under namespaced prefixes.
    \item \textbf{RabbitMQ} ensures reliable command/event delivery and provides natural back-pressure.
    \item \textbf{Browser storage} (IndexedDB/sessionStorage) buffers local chunks to survive transient network hiccups during recording.
\end{itemize}

\section{Event Contracts}

\begin{table}[H]
\centering
\resizebox{\textwidth}{!}{%
\begin{tabular}{>{\bfseries}p{3.8cm}p{4.6cm}p{6.0cm}p{4.6cm}}
\toprule
Event & Producer & Payload Highlights & Primary Consumers \\
\midrule
\texttt{session.created} &
Media Recording Service &
Session ID, room code, host ID, timestamps &
Analytics, audit trail, notification systems \\
\texttt{chunk.uploaded} &
Media Recording Service &
Chunk ID, recording ID, sequence number, MinIO object, payload size &
Progress dashboards, retry monitors \\
\texttt{upload.completed} &
Media Recording Service &
Recording ID, total chunks, total bytes, participant ID &
Processing orchestration, email/SMS hooks \\
\texttt{recording.processed} &
Media Processing Worker &
Session ID, recording ID, track type, processed MinIO object, file size &
Frontend polling, transcription pipeline, archival service \\
\bottomrule
\end{tabular}}
\caption{Event-driven contract summary across services.}
\label{tab:events}
\end{table}

\clearpage
\section{Deployment Topology}

\begin{figure}[H]
\centering
\resizebox{0.9\textwidth}{!}{%
\begin{tikzpicture}[
    box/.style={rectangle, draw=black, rounded corners, fill=blue!10, minimum width=4.0cm, minimum height=1.0cm, align=center},
    infra/.style={rectangle, draw=black, fill=orange!15, minimum width=3.8cm, minimum height=1.0cm, align=center},
    arrow/.style={-{Latex[length=2.5mm]}, thick}
]
    \node[box] (frontend) {podcast-frontend\\(Next.js 14)};
    \node[box, below=0.9cm of frontend] (recording) {media-recording-service\\(FastAPI)};
    \node[box, below=0.9cm of recording] (processing) {media-processing-service\\(FastAPI)};
    \node[box, below=0.9cm of processing] (worker) {media-processing-worker\\(FFmpeg)};

    \node[infra, right=4.0cm of recording] (rabbit) {RabbitMQ};
    \node[infra, above=1.2cm of rabbit] (minio) {MinIO};
    \node[infra, below=1.2cm of rabbit] (postgres) {PostgreSQL};
    \node[infra, below=1.0cm of postgres] (redis) {Redis};

    \draw[arrow] (frontend) -- (recording);
    \draw[arrow] (recording) -- (rabbit);
    \draw[arrow] (recording) -- (minio);
    \draw[arrow] (recording) -- (postgres);
    \draw[arrow] (processing) -- (rabbit);
    \draw[arrow] (worker) -- (rabbit);
    \draw[arrow] (worker) -- (minio);
\end{tikzpicture}}
\caption{Docker Compose deployment topology with shared infrastructure.}
\label{fig:deployment}
\end{figure}

Operational considerations include container health checks, environment overrides for internal hostnames (e.g., \texttt{rabbitmq:5672}), and the ability to scale the worker horizontally for additional throughput.

\section{Repository Map}

\begin{longtable}{>{\bfseries}p{4.0cm}p{10cm}}
\caption{Source tree responsibilities.}
\label{tab:code}\\
\toprule
Directory / File & Responsibility \\
\midrule
\endfirsthead
\toprule
Directory / File & Responsibility \\
\midrule
\endhead
\bottomrule
\endfoot
media-recording-service/src/domain/models/ & Aggregate definitions for recordings, chunks, and uploads. \\
media-recording-service/src/application/services/ & Use cases coordinating domain logic and events. \\
media-recording-service/src/adapters/inbound/ & REST routes, WebSocket signalling, simplified HTTP upload entry point. \\
media-recording-service/src/adapters/outbound/ & MinIO storage adapter, RabbitMQ publisher, repository implementations. \\
media-recording-service/src/processors/media\_processing\_worker.py & Background worker executing FFmpeg stitching. \\
media-processing-service/src/adapters/inbound/rest/ & Administrative endpoints and job orchestration APIs. \\
podcast-frontend/src/app/ & Next.js pages implementing host/guest flows. \\
podcast-frontend/src/hooks/ & Custom hooks for WebRTC, recording management, and upload orchestration. \\
docker-compose.yml & Service topology, health checks, environment overrides, volume definitions. \\
postman/*.json & Collections for exercising recording and processing APIs. \\
\end{longtable}

\section{Testing \& Validation}

\begin{table}[H]
\centering
\resizebox{\textwidth}{!}{%
\begin{tabular}{>{\bfseries}p{3.8cm}p{4.2cm}p{5.8cm}p{3.0cm}}
\toprule
Test Category & Tools / Scripts & Coverage Highlights & Status \\
\midrule
Static/Syntax & \texttt{python -m compileall}, \texttt{npm run build} & FastAPI modules, worker scripts, Next.js pages compile without error & Pass \\
API Flows & Postman collections & Session creation/join, recording start/stop, chunk upload, processing enqueue, job status & Pass \\
Manual Scenario & Browser walkthrough & WebRTC handshake, chunk upload verification in MinIO, worker output, processed download & Pass \\
Fault Injection & Manual (network toggle) & Verified queued commands survive restarts; partial uploads retried successfully & Pass \\
\bottomrule
\end{tabular}}
\caption{Testing matrix summarising validation activities.}
\label{tab:tests}
\end{table}

\section{Scenario Walkthrough}\label{sec:scenario}

\begin{enumerate}[leftmargin=*]
    \item \textbf{Create session}: Host calls \texttt{/api/sessions/create}; recording service emits \texttt{session.created} and returns room code.
    \item \textbf{Join \& handshake}: Guest enters room code, WebSocket signalling completes SDP exchange, and WebRTC media streams are established.
    \item \textbf{Record \& upload}: Every five seconds browsers post chunks to \texttt{/api/uploads/chunk}; service stores metadata, streams bytes to MinIO, and emits \texttt{chunk.uploaded}.
    \item \textbf{Enqueue processing}: Host stops recording; frontend calls \texttt{/api/uploads/recording/\{id\}/enqueue-processing}. A command is pushed to RabbitMQ.
    \item \textbf{Worker stitch}: Worker downloads chunks, runs FFmpeg concat, uploads processed asset, and emits \texttt{recording.processed}.
    \item \textbf{Delivery}: Frontend polls progress endpoints and presents download links; downstream services can subscribe to processed events.
\end{enumerate}

\section{Quality Attribute Discussion}

\begin{table}[H]
\centering
\resizebox{\textwidth}{!}{%
\begin{tabular}{>{\bfseries}p{3.6cm}p{5.8cm}p{6.0cm}}
\toprule
Attribute & Design Mechanisms & Observations / Trade-offs \\
\midrule
Reliability & Checksum validation, durable queues, idempotent object keys & Protects against chunk corruption and ensures recovery from worker restarts. \\
Scalability & Stateless APIs, horizontally scalable workers, RabbitMQ back-pressure & Worker replicas can be added without touching the recording service. \\
Extensibility & Ports/adapters and domain events & Storage broker swaps and new event consumers (transcription) do not affect core logic. \\
Maintainability & Layered code structure, repository mapping & Developers navigate by domain or adapter type; unit tests can target ports. \\
Observability & Health endpoints, RabbitMQ/MinIO consoles, structured logging & Baseline observability in place; ready for Prometheus exporters. \\
Security (future work) & Authentication, TLS termination, secrets management & Deferred for course scope; design supports integration with auth gateway and vault. \\
\bottomrule
\end{tabular}}
\caption{Quality attributes and supporting design choices.}
\label{tab:quality}
\end{table}

\section{Future Enhancements}

\begin{itemize}[leftmargin=*]
    \item Implement PostgreSQL adapters for persistent repositories and replace in-memory stores.
    \item Add authentication/authorisation (JWT/OAuth) to secure recording endpoints.
    \item Introduce observability stack (Prometheus, Grafana, Loki) for metrics and log aggregation.
    \item Extend processing workflows for loudness normalisation, automated transcription, and publishing.
    \item Automate resilience tests (chaos experiments) to validate queue replay and retry logic.
\end{itemize}

\section{Lessons Learned}

\begin{enumerate}[leftmargin=*]
    \item Hexagonal architecture kept business rules infrastructure-agnostic, allowing the worker to reuse the same domain layer.
    \item Event-driven coordination prevented long-running FFmpeg operations from blocking user-facing APIs.
    \item Chunk-first recording provided resilience to network instability and enabled deterministic reprocessing.
    \item Containerisation delivered a one-command demo experience for faculty and peers.
\end{enumerate}

\section{Conclusion}

PodcastHub demonstrates how graduate-level architectural concepts translate into a functional, media-centric platform. By decomposing responsibilities into bounded contexts, orchestrating services through events, and enforcing hexagonal boundaries, the system remains reliable, scalable, and extensible. The accompanying diagrams and tables provide a clear narrative for both professors and students reviewing the project.
\end{document}
